\item \textbf{Problema de repaso.}
Considere el aparato que se muestra en la figura 17.15 y se describe en el ejemplo 17.4. Supongamos que el número de antinodos en la figura es un valor arbitrario $n$.
\begin{enumerate}
    \item Encuentre una expresión para el radio de la esfera en el agua como una función solo de $n$.
    \item ¿Cuál es el mínimo valor permitido de $n$ para una esfera de tamaño distinto de cero?
    \item ¿Cuál es el radio de la esfera más grande que se producirá en una onda en la cuerda?
    \item ¿Qué pasa si se utiliza una esfera más grande?
\end{enumerate}